\chapter{Background and Literature Review} \label{chapter:background}

\section{Defining Material Interfaces} \label{section:material_interfaces}

This subsection establishes a foundational understanding of material interfaces, the boundary regions where two distinct phases, often different materials, meet. These interfaces are critical in determining the structural, electronic, and thermal behaviour of composite systems. In solid-state devices, subtle atomic arrangements and local electronic environments at interfaces can profoundly influence macroscopic properties. This is especially true in modern micro- and nanoscale technologies, where the continual miniaturisation of components has shifted the paradigm from bulk-dominated behaviour to interface-dominated performance. As devices shrink, interfaces no longer represent isolated regions but instead constitute a significant fraction of the material, making their precise definition and understanding essential for predictive modelling and materials design .

\subsection{What is a material interface in the context of this research?}

A material interface is the region where two distinct solids meet and form a junction, commonly arising in heterostructures, grain boundaries, or thin film systems. These regions are not only geometrical boundaries but exhibit unique atomic, electronic, and thermal properties that distinguish them from the bulk phases of the parent materials. The local structure at the interface can include strain fields, atomic reconstruction, and intermixing of constituent elements, all of which can significantly influence the behaviour of the composite system.

In this research, material interfaces are studied predominantly through \textit{ab initio} methods, particularly Density Functional Theory (DFT), with a focus on predicting and analysing the structural and energetic properties of 2D|2D and 2D|3D junctions. Interfaces are modelled as systems with finite thickness that bridge two crystalline lattices. Depending on their geometric orientation, interfaces are categorised as either \emph{stacked}, where a 2D layer rests atop a substrate, or \emph{lateral}, where edge-on growth leads to in-plane bonding between dissimilar materials.
Several methodologies have been developed to address the challenges of constructing realistic interface models. The ARTEMIS tool enables high-throughput generation of interface structures through lattice matching and surface termination selection, allowing systematic exploration of the complex interfacial phase space. Similarly, the RAFFLE method uses statistical learning and void-filling algorithms to construct interface phases that are energetically favourable and representative of plausible morphologies .

Thus, in the context of this research, a material interface is not merely a geometric boundary but an emergent, functionally distinct region. It governs critical physical properties, demands precise structural modelling, and serves as a fertile ground for material innovation and device engineering.

\subsubsection{Types of interfaces}

Material interfaces can be broadly categorised by their crystallographic orientation, dimensionality, and degree of lattice registry. These classifications are instrumental in understanding interfacial phenomena and in designing simulation protocols that are physically realistic and computationally tractable.

\paragraph{By geometry: stacked vs lateral.} Stacked interfaces occur when one material is deposited atop another, forming an out-of-plane junction. This is characteristic of many van der Waals heterostructures, where weak interlayer forces enable a range of stacking configurations and rotational alignments. In contrast, lateral interfaces arise when two materials are grown side-by-side, often leading to in-plane covalent bonding at the junction. These interfaces are especially relevant for in-plane heterostructures such as MoS$_2$|WS$_2$ or graphene-based edge-junctions.

\paragraph{By coherence: coherent, semi-coherent, incoherent.} Interfaces can also be distinguished by the degree of lattice match across the boundary:

\begin{itemize}
    \item \textbf{Coherent interfaces} exhibit perfect registry between atomic planes, typically occurring when lattice mismatch is negligible. These interfaces preserve the translational symmetry across the junction and are often found in homoepitaxial or pseudomorphic systems.
    \item \textbf{Semi-coherent interfaces} accommodate lattice mismatch through periodic misfit dislocations. These dislocations relieve interfacial strain while maintaining partial registry, allowing larger domains of coherency interspersed with localised defects.
    \item \textbf{Incoherent interfaces} display no periodicity match across the interface, leading to structurally disordered junctions. These are often found in polycrystalline or amorphous/crystalline composites and can act as strong phonon or electron scattering centres.
\end{itemize}

\paragraph{By chemical contrast: heterojunctions vs homojunctions.} Chemical composition also plays a defining role. \emph{Homojunctions} involve structurally similar regions with subtle variations (e.g. doping or strain-modulated domains), while \emph{Heterojunctions} involve distinct materials with differing electronegativity, band gaps, and ionic/covalent character. Heterojunctions are central to semiconductor devices, solar cells, and catalysis, often giving rise to emergent interfacial dipoles, states, or reconstructed bonding motifs.

These classifications are not mutually exclusive; real interfaces may be simultaneously lateral and semi-coherent, or stacked but chemically graded. Understanding these distinctions is essential when generating or interpreting simulation results for both idealised and realistic interfacial systems.

\subsection{What physical or electronic properties are commonly affected by interfaces?}

Material interfaces influence a wide array of physical and electronic properties, often determining the functional behaviour of composite systems and devices. These effects are particularly pronounced in systems involving 2D materials, complex oxides, or strongly correlated electron systems.

\subsubsection{Band Alignment and Electronic Transport}

Perhaps the most critical electronic property influenced by interfaces is the band alignment between materials. Misalignment of conduction and valence band edges creates energy barriers or wells that determine whether an interface behaves as a Schottky barrier, Ohmic contact, or heterojunction. These alignments are crucial for device operation in field-effect transistors, diodes, and photovoltaics. Interfacial dipoles, charge transfer, and electronic reconstruction can shift the effective band edges, often invalidating simple rules like Anderson\rqs model.

\subsubsection{Charge Carrier Dynamics}

Interfaces often host localised states arising from atomic mismatch, bonding reconstruction, or impurity segregation. These can act as traps or recombination centres, significantly affecting carrier lifetimes and mobility. In 2D|3D systems, tunnelling, leakage, or enhanced recombination may occur depending on the barrier height and band offset at the junction.

\subsubsection{Dielectric Response and Permittivity}

The dielectric properties of materials can be strongly altered at interfaces. In systems like \ce{CaCu3Ti4O12}, colossal permittivity arises not from the bulk crystal but from insulating grain boundaries and a semiconducting interior, exemplifying the Maxwell-Wagner effect. Interface-induced polarisation, charge accumulation, and ionic displacements all contribute to enhanced dielectric responses.

\subsubsection{Thermal Conductivity}

Thermal transport across interfaces is typically suppressed due to phonon scattering, mass disorder, and acoustic impedance mismatch. This results in thermal boundary resistance (or Kapitza resistance), a key factor in thermoelectric design and heat management in microelectronics. In 2D heterostructures, the interface often acts as a dominant phonon scattering site, reducing the effective thermal conductivity below that of either constituent layer.

\subsubsection{Structural Stability and Phase Formation}

Interfaces can stabilise new or metastable phases not accessible in bulk. For instance, the interface between \ce{BaTiO3} and \ce{SiO2} has been shown to promote the spontaneous formation of the fresnoite phase \ce{Ba2TiSi2O8}, altering both electronic and mechanical response. A more nuanced example is the graphene|\ce{MgO} interface: while monolayer MgO confined between graphene layers often adopts a rocksalt-like structure, this is not simply a case of one phase being energetically favoured. Instead, the local environment, including charge transfer and encapsulation, stabilises the rocksalt phase relative to the expected hexagonal phase in unsupported monolayers. However, experimental and theoretical studies suggest that a mixed or hybrid phase may also emerge with phase stability dependent on thickness and interface strain.

\subsubsection{Mechanical Properties and Elastic Moduli}

Mechanical properties such as bulk modulus, shear strength, and adhesion energy are modified at interfaces due to altered bonding environments and lattice mismatch. These variations are important in the design of composite materials, flexible electronics, and nanoindentation studies.

\subsubsection{Optical Properties and Absorption}

Interfaces affect the optical absorption spectrum by introducing new electronic states and modifying the joint density of states. This can lead to bandgap narrowing or broadening, shifts in refractive index, and enhancement of excitonic effects; features exploited in photodetectors and photovoltaic devices.

\paragraph{In Summary} material interfaces act as multifunctional domains that modulate a broad spectrum of material properties. Their control and prediction are vital for the design of next-generation electronic, optoelectronic, and thermal devices.